\begingroup
\let\clearpage\relax
\let\cleardoublepage\relax
\chapter*{Preface}
\endgroup
\addcontentsline{toc}{chapter}{Preface}

\vspace{1.5em}
\begin{center}
\begin{minipage}{0.62\textwidth}
\centering
\itshape
How easy it is\par
to look out across the sea\par
while you rest on the white sand.\par\medskip
Perhaps you never noticed,\par
but you swam so swiftly\par
that no one saw you leave.\par\medskip
The waves carried you\par
while you simply lay there.\par
You reached for air\par
as your vision faded.\par\medskip
You gasp to call out\par
and that was when your chest burned.
\end{minipage}
\end{center}
\vspace{2em}

When I came to UiO, Morten read my background and said: ``You know PINNs, you know
many-body quantum systems. I want you to do PINNs on quantum dots.'' I said
``sure.'' I knew enough about each of those things separately to imagine that putting
them together would mostly be a technical problem.

It was not.

The difficult part was rarely producing a result. The difficult part was knowing what
a result meant. When something stopped improving, there were always several possible
explanations. Perhaps the model was wrong for the problem, perhaps the optimisation was
failing, perhaps the physics itself had become more difficult, or perhaps I had simply
made a mistake somewhere along the way. Much of this thesis grew out of trying to tell
those possibilities apart.

The trouble with having so many explanations was that, as one after another failed to
account for what I was seeing, a much simpler one began to acquire weight: perhaps it
simply could not be done. At first it was only another possibility. Gradually, it became
difficult to see past it.

For a while, I stopped expecting the project to succeed.

I kept working, but for a different reason. If it could not be done, I wanted at least
to understand why. In some ways, that became a more difficult problem than the original
one. The worst periods were not those in which the method failed, but those in which I
knew that it had failed and could not say what the failure meant.

I was fortunate to spend part of that period in Rome. 
The city’s ethereal quality lifted my spirits and lent the days a sense of something 
larger than the problem I happened to be working on.
Yet it was my girlfriend Elida, more than anything, 
who gave me a steadiness I could not have supplied for myself. 
Her patience, warmth, 
and faith in me carried me through more of this work than I know how to put into words.

Eventually, some of the mystery began to give way. Not every obstacle turned out to be
profound. Some were almost embarrassingly ordinary. Things I had taken for limitations
of the method occasionally turned out to be limitations of my own implementation; other
failures, once understood properly, pointed toward more interesting questions than the
ones I had originally been asking.

That changed the way I looked at the project. A disappointing result was no longer
merely something to overcome. Sometimes it contained information about the model,
sometimes about the physics, and sometimes about the question itself. What began as an
attempt to obtain a sufficiently good answer gradually became an attempt to understand
what the answers, and the failures, were telling me.

\bigskip

I owe a warm and sincere thanks to my family, who have always made sure there was a home
for me to return to, whenever I needed one. And to my friends, for countless
conversations in which we tried to work out what we actually understood, and usually
discovered that it was less than we thought.

To Morten, again, and properly: thank you for the assignment, for the two years we have
had together, and above all for a genuine and undisguised interest in what I was finding
at every stage, including the stages where what I found was that I had been wrong. That
interest is most of what encouragement actually consists of.

\bigskip

When I look back, it is easy to make a catalogue of what I would now do differently.
Most of that knowledge, however, was purchased by the failures and could not have been
read off in advance. 

I began this project thinking that the difficult part would be finding an answer to the
many-body problem. More often, it was discovering which question a disappointing result
was actually answering.

\clearpage
